% Gerado por src/python/thesis/build_tab_engenharia.py -- nao editar a mao.
% Regenerar com: make thesis-tables
% Fonte: results/engineering_suite/engineering_suite_summary_main.csv
% Fonte: results/engineering_suite/engineering_suite_pairwise_main.csv
\begin{table}[!htb]
    \centering
    \small
    \caption{Desempenho do IVF/SPEA2 em problemas de engenharia com restrições: comparação com o hospedeiro e com os demais algoritmos.}
    \label{tab:engenharia_wtl}
    \begin{tabular*}{\textwidth}{@{\extracolsep{\fill}}lccrrcc@{}}
        \toprule
        \multicolumn{3}{@{}c}{} & \multicolumn{3}{c}{\textbf{Contra o SPEA2}} & \multicolumn{1}{c@{}}{\textbf{Contra todos}} \\
        \cmidrule(lr){4-6}\cmidrule(l){7-7}
        \textbf{Problema} & \textbf{$M$} & \textbf{Métrica} & \textbf{$\Delta$ (\%)} & \textbf{$p$} & \textbf{Resultado} & \textbf{V/E/D} \\
        \midrule
        \multirow{2}{*}{RWMOP9} & \multirow{2}{*}{2} & IGD & $+1{,}17$ & $0{,}0018$ & Vitória & 8/0/0 \\
         &  & HV & $+0{,}01$ & $0{,}0035$ & Vitória & 5/0/3 \\
        \midrule
        \multirow{2}{*}{RWMOP21} & \multirow{2}{*}{2} & IGD & $+0{,}27$ & $0{,}8604$ & Empate & 3/2/3 \\
         &  & HV & $0{,}00$ & $0{,}1749$ & Empate & 3/2/3 \\
        \midrule
        \multirow{2}{*}{RWMOP8} & \multirow{2}{*}{3} & IGD & $+4{,}96$ & $0{,}8563$ & Empate & 1/1/5 \\
         &  & HV & $+0{,}31$ & $0{,}5444$ & Empate & 1/1/5 \\
        \bottomrule
    \end{tabular*}
    \par\smallskip
    \parbox{\textwidth}{\footnotesize
    Contra o SPEA2: $\Delta$ é a diferença relativa entre as medianas, positiva quando favorece o IVF/SPEA2, e $p$ vem do teste de soma de postos de Wilcoxon bicaudal sobre 60 execuções de cada algoritmo, com $\alpha = 0{,}05$; com correção de Holm entre os três problemas de cada métrica, nenhuma decisão muda. Contra todos: V/E/D do IVF/SPEA2 frente a cada um dos 8 algoritmos restantes, incluído o SPEA2, pelo mesmo teste. No RWMOP8, a cobertura de execuções válidas é heterogênea (SPEA2+SDE tem 18 de 60; MOEA/D, nenhuma), e a contagem soma 7 comparadores.
    }
\end{table}
